\section{Generalization to Higher Spins}
\label{sec:HigherSpins}
To systematically address the extension to higher spins, we adopt the Bargmann-Wigner (BW) formalism,
which provides a rigorous group-theoretical framework for constructing relativistic wave equations
for particles of arbitrary spin $s$. Established by V. Bargmann and E. P. Wigner \cite{bargmann1997a},
this method posits that the wavefunction of a particle with spin $s$ can be represented by a multispinor
field $\Psi_{\alpha_1 \alpha_2 \dots \alpha_{2s}}$ of rank $2s$, which effectively describes a composite
state of $2s$ fundamental spin-1/2 fermions. To ensure that this object transforms irreducibly under
the Poincar\'{e} group, the multispinor must satisfy two conditions: it must be totally symmetric 
under the exchange of any pair of spinor indices, and it must satisfy the Dirac equation for each 
index independently:
\begin{equation}
    \left( i\,\gamma^\mu \partial_\mu - m \right)_{\alpha_k \beta}\,
    \Psi_{\alpha_1 \dots \beta \dots \alpha_{2s}}(x)
    = 0,
    \qquad \text{for } k = 1, \dots, 2s .
\end{equation}

To address the specific extension to vector bosons, we adopt the methodology established by Watson in their
generalization of the Proca equation \cite{Watson2022}. Their approach serves as the guiding framework 
for this analysis. Specifically, for the vector boson case ($s=1$), the formalism begins with the 
Bargmann-Wigner equations for a symmetric rank-2 spinor $\Psi_{\alpha\beta}$. Crucially, the standard
scalar mass $m$ is replaced by the chiral mass operator $\hat{M}$ derived from the fermionic theory, 
applied to each spinor index independently:
\begin{equation}
    (i \gamma^\mu \partial_\mu - \hat{M})_{\alpha \lambda} \Psi_{\lambda \beta}(x) = 0, \quad \text{with} \quad \hat{M} = m e^{-i\alpha\gamma^5}.
\end{equation}

The wavefunction is then expanded in terms of the Clifford algebra generators. The symmetry constraint 
intrinsic to the Bargmann-Wigner formalism for $s=1$ requires the spinor $\Psi_{\alpha\beta}$ to be symmetric,
which effectively eliminates the scalar, pseudoscalar, and axial-vector components. Consequently, the wavefunction
reduces naturally to a superposition of only the vector ($A_\mu$) and tensor ($F_{\mu\nu}$) sectors:
\begin{equation}
    \Psi(x) = \left( \gamma^\mu A_\mu(x) + \frac{1}{2}\sigma^{\mu\nu} F_{\mu\nu}(x) \right) C,
\end{equation}
where $C$ is the charge conjugation matrix. Substituting this ansatz back into the fundamental chiral equations
leads to a decoupling of the components. While the standard theory yields the Proca equation $\partial_\mu F^{\mu\nu} 
+ m^2 A^\nu = 0$, the chiral extension results in the Generalized Proca Equation:
\begin{equation}
    \partial_\mu F^{\mu\nu} + m^2 \cos^2(\alpha) A^\nu + \dots = 0.
\end{equation}

It is important to clarify that the appearance of the $\cos^2(\alpha)$ factor in this sector does not 
violate the relativistic structure of the theory; rather, it acts as a scaling parameter that defines
an effective physical mass $M_{\text{eff}} = m \cos(\alpha)$. This implies that the chiral angle can
modulate the observable mass of the boson without altering the underlying geometric consistency of 
the Poincar\'{e} group.

\subsection{Extension to the Rarita-Schwinger Field ($s=3/2$)}

Building on this success, we apply the procedure to the Rarita-Schwinger field ($s=3/2$). 
To maintain full consistency with the underlying Bargmann-Wigner theory, we treat the field 
strictly as a totally symmetric rank-3 spinor $\Psi_{\alpha\beta\gamma}$. In this formalism, 
the field is a composite state of three fundamental spin-1/2 fermions.

The chiral mass generation mechanism is applied to each of the three constituent spinor indices. 
Since the chiral rotation $e^{-i\alpha\gamma^5}$ induces a projection factor of $\cos(\alpha)$ 
for a single spinor (as seen in the Dirac case), the composite mass term for a rank-3 object 
must scale as the product of the projections of its constituents. This leads to a cubic scaling law:
\begin{equation}
    \hat{M}_{\text{composite}} \rightarrow m \cos^3(\alpha).
\end{equation}

Consequently, the modified Chiral Rarita-Schwinger equation takes the form:
\begin{equation}
    \epsilon^{\mu\nu\rho\sigma} \gamma^5 \gamma_\nu \partial_\rho \psi_\sigma + i m \cos^3(\alpha) \sigma^{\mu\nu} \psi_\nu = 0.
\end{equation}

This result suggests a general scaling law for the effective mass of a particle with spin $s$:
$M_{\text{eff}}^{(s)} \approx m \cos^{2s}(\alpha)$. This theoretical prediction aligns with the 
linear scaling found for Dirac fermions ($2s=1 \rightarrow \cos^1$) and the quadratic scaling found 
for Proca bosons ($2s=2 \rightarrow \cos^2$), offering a unified geometric interpretation of mass modulation across the spin spectrum.

Finally, we note that for higher-spin fields, the introduction of non-trivial mass matrices requires careful checking of the constraint algebra ($\gamma^\mu\psi_\mu=0$ and $\partial^\mu\psi_\mu=0$) to avoid the propagation of acausal modes. This is a known issue in massive Rarita-Schwinger theories referred to as the Velo-Zwanziger problem \cite{VeloZwanziger}. While a full causality analysis is beyond the scope of this proposal, future work must verify the preservation of these constraints under chiral mass modifications.

\subsection{Discussion on Generalization via Projection Operators}

The explicit construction of equations for higher spins (like $s=3/2$ or $s=2$) using standard component 
notation becomes exponentially complex. To solve this, Watson and Musielak propose a powerful generalization 
based on Orthogonal Idempotents.

Instead of manually constructing the Lagrangians, this method utilizes projection operators $\hat{P}$ to 
define the physical subspaces directly. For a system of $N$ constituent spinors (which would form a 
spin-$N/2$ particle), the generalized projection operator is defined as the tensor product of the individual 
helicity projectors:
\begin{equation}
    \hat{P}_{s_1, s_2, \dots, s_N}(\hat{n}) = \hat{P}_{s_1}(\hat{n}) \otimes \hat{P}_{s_2}(\hat{n}) \otimes \dots \otimes \hat{P}_{s_N}(\hat{n}),
\end{equation}
where each $\hat{P}_{s_i}(\hat{n})$ projects the spin state relative to a quantization axis $\hat{n}$. For 
the specific case of the Rarita-Schwinger field ($N=3$), one selects the symmetric subspace of three entangled spinors.

This projector formalism is superior because it automatically handles the chiral constraints. By defining the 
chiral basis for each projector (parameterized by angles $\alpha_1, \alpha_2, \dots$), the resulting equation 
of motion emerges naturally with the mass terms already modulated by the relative chiral angles. Thus, the 
effective mass $M_{\text{eff}}$ is not an ad-hoc insertion but an eigenvalue of the generalized projection 
operator acting on the vacuum state:
\begin{equation}
    \hat{P}_{\text{total}} \Psi = M_{\text{eff}} \Psi.
\end{equation}
This confirms that the chiral degrees of freedom are fundamental to the algebraic structure of the 
Poincar\'{e} group representations, providing a rigorous path to derive wave equations for any arbitrary spin $s$.
